\documentclass[aspectratio=169,11pt]{beamer}

\usetheme{Madrid}
\usecolortheme{seahorse}
\setbeamertemplate{navigation symbols}{}
\setbeamertemplate{footline}[frame number]

\usepackage{fontspec}
\setsansfont{DejaVu Sans}
\setmonofont{DejaVu Sans Mono}[Scale=0.88]
\usepackage{booktabs}
\usepackage{tabularx}
\usepackage{array}
\usepackage{xcolor}
\usepackage{listings}
\usepackage{ragged2e}
\usepackage{tikz}
\usetikzlibrary{positioning,arrows.meta}

\definecolor{NEUBlue}{HTML}{1F4E79}
\definecolor{NEULight}{HTML}{EAF3FF}
\definecolor{SQLBg}{HTML}{F8FAFC}
\definecolor{SQLKey}{HTML}{1D4ED8}
\definecolor{SQLComment}{HTML}{64748B}
\definecolor{SQLString}{HTML}{047857}
\setbeamercolor{structure}{fg=NEUBlue}
\setbeamercolor{frametitle}{fg=white,bg=NEUBlue}
\setbeamercolor{title}{fg=white,bg=NEUBlue}
\setbeamerfont{frametitle}{size=\Large,series=\bfseries}
\setbeamerfont{title}{size=\LARGE,series=\bfseries}

\lstdefinelanguage{SQL}{
  morekeywords={SELECT,FROM,WHERE,INNER,JOIN,LEFT,RIGHT,ON,AS,ORDER,BY,DESC,ASC,LIMIT,IS,NULL,AND,OR,NOT,IN,BETWEEN,LIKE,USE,SHOW,TABLES,DESC,DESCRIBE,COUNT,SUM,AVG,MIN,MAX,GROUP,HAVING},
  sensitive=false,
  morecomment=[l]{--},
  morestring=[b]',
}
\lstset{
  language=SQL,
  basicstyle=\ttfamily\scriptsize,
  keywordstyle=\color{SQLKey}\bfseries,
  commentstyle=\color{SQLComment},
  stringstyle=\color{SQLString},
  backgroundcolor=\color{SQLBg},
  frame=single,
  rulecolor=\color{black!15},
  breaklines=true,
  columns=fullflexible,
  keepspaces=true,
  showstringspaces=false,
  xleftmargin=0.5em,
  xrightmargin=0.5em,
  aboveskip=0.6em,
  belowskip=0.4em
}

\newcommand{\tbl}[1]{\texttt{#1}}
\newcommand{\col}[1]{\texttt{#1}}
\newcommand{\kw}[1]{\textbf{\texttt{#1}}}
\newcommand{\hint}[1]{\begin{block}{Ghi nhớ}#1\end{block}}
\newcommand{\sourceurl}{\footnotesize Nguồn: \texttt{vudmvn.github.io/dbms-course/.../select-statement-2}}

\title[SELECT và JOIN]{Tutorial: Sử dụng \texttt{SELECT} và \texttt{JOIN} với \texttt{classicmodels}}
\subtitle{MySQL Querying Data -- Select Statement 2}
\author{DBMS Course}
\institute{Smart Logistics and Supply Chain Management Lab}
\date{\today}

\begin{document}

% 1
\begin{frame}
  \titlepage
  \vspace{-0.7em}
  \begin{center}
    \sourceurl
  \end{center}
\end{frame}

% 2
\begin{frame}{Mục tiêu học tập}
\begin{columns}[T,onlytextwidth]
\column{0.48\textwidth}
\begin{block}{Nhắc lại truy vấn một bảng}
\begin{itemize}
  \item Viết \kw{SELECT ... FROM ...}
  \item Lọc bằng \kw{WHERE}
  \item Sắp xếp bằng \kw{ORDER BY}
  \item Giới hạn kết quả bằng \kw{LIMIT}
\end{itemize}
\end{block}
\column{0.48\textwidth}
\begin{block}{Kết hợp dữ liệu nhiều bảng}
\begin{itemize}
  \item Hiểu khóa chính -- khóa ngoại
  \item Viết \kw{INNER JOIN}, \kw{LEFT JOIN}
  \item Nối 3, 4, 5 bảng
  \item Viết \kw{SELF JOIN} đơn giản
\end{itemize}
\end{block}
\end{columns}
\end{frame}

% 3
\begin{frame}{CSDL mẫu \texttt{classicmodels}}
\begin{tabularx}{\textwidth}{>{\ttfamily}l X >{\ttfamily}l X}
\toprule
productlines & Dòng sản phẩm & products & Thông tin sản phẩm \\offices & Văn phòng & employees & Nhân viên \\customers & Khách hàng & payments & Thanh toán \\orders & Đơn hàng & orderdetails & Chi tiết đơn hàng \\
\bottomrule
\end{tabularx}
\vspace{0.8em}
\hint{Trong bài này, \texttt{classicmodels} được dùng để luyện cách đọc dữ liệu từ nhiều bảng có quan hệ với nhau.}
\end{frame}

% 4
\begin{frame}[fragile]{Chuẩn bị môi trường}
\begin{lstlisting}
USE classicmodels;
SHOW TABLES;
DESC products;
DESC customers;
DESC orderdetails;
\end{lstlisting}
\begin{itemize}
  \item \kw{USE}: chọn cơ sở dữ liệu đang làm việc.
  \item \kw{SHOW TABLES}: xem danh sách bảng.
  \item \kw{DESC}: xem cấu trúc bảng, tên cột, kiểu dữ liệu, khóa.
\end{itemize}
\end{frame}

% 5
\begin{frame}{Các quan hệ chính cần nhớ}
\scriptsize
\begin{tabularx}{\textwidth}{l l l l}
\toprule
Bảng con & Khóa ngoại & Bảng cha & Khóa chính \\
\midrule
products & productLine & productlines & productLine \\
employees & officeCode & offices & officeCode \\
employees & reportsTo & employees & employeeNumber \\
customers & salesRepEmployeeNumber & employees & employeeNumber \\
payments & customerNumber & customers & customerNumber \\
orders & customerNumber & customers & customerNumber \\
orderdetails & orderNumber & orders & orderNumber \\
orderdetails & productCode & products & productCode \\
\bottomrule
\end{tabularx}
\vspace{0.5em}
\centering
\texttt{customers -> orders -> orderdetails -> products -> productlines}
\end{frame}

% 6
\begin{frame}{Tư duy khi viết JOIN}
\begin{enumerate}
  \item Muốn lấy thông tin gì?
  \item Các cột đó nằm ở bảng nào?
  \item Cặp khóa nào nối hai bảng?
  \item Cần giữ dòng không khớp hay chỉ lấy dòng khớp?
\end{enumerate}
\vspace{0.7em}
\begin{block}{Công thức thực hành}
\centering
\texttt{SELECT columns -> FROM main\_table -> JOIN related\_table -> ON key\_pair}
\end{block}
\end{frame}

% 7
\begin{frame}[fragile]{Nhắc lại \texttt{SELECT} cơ bản}
\begin{lstlisting}
SELECT customerNumber, customerName, city, country
FROM customers;
\end{lstlisting}
\begin{lstlisting}
SELECT productCode, productName, productLine, MSRP
FROM products;
\end{lstlisting}
\begin{itemize}
  \item Chỉ chọn các cột cần dùng.
  \item Khi nhiều bảng có cùng tên cột, cần ghi rõ \texttt{table.column} hoặc \texttt{alias.column}.
\end{itemize}
\end{frame}

% 8
\begin{frame}[fragile]{Lọc dữ liệu bằng \texttt{WHERE}}
\begin{lstlisting}
SELECT customerNumber, customerName, city, country
FROM customers
WHERE country = 'USA';
\end{lstlisting}
\begin{lstlisting}
SELECT orderNumber, orderDate, status
FROM orders
WHERE status = 'Cancelled';
\end{lstlisting}
\hint{\kw{WHERE} lọc dòng. Với \texttt{JOIN}, có thể lọc theo cột của bất kỳ bảng nào đã nối.}
\end{frame}

% 9
\begin{frame}[fragile]{Sắp xếp và giới hạn kết quả}
\begin{lstlisting}
SELECT productCode, productName, MSRP
FROM products
ORDER BY MSRP DESC
LIMIT 5;
\end{lstlisting}
\begin{itemize}
  \item \kw{ORDER BY}: sắp xếp kết quả.
  \item \kw{DESC}: giảm dần; \kw{ASC}: tăng dần.
  \item \kw{LIMIT}: giới hạn số dòng trả về.
\end{itemize}
\end{frame}

% 10
\begin{frame}{Vì sao cần \texttt{JOIN}?}
\begin{columns}[T,onlytextwidth]
\column{0.48\textwidth}
\begin{block}{Dữ liệu được tách bảng}
\begin{itemize}
  \item \tbl{orders} chỉ lưu \col{customerNumber}.
  \item Tên khách hàng nằm trong \tbl{customers}.
  \item Chi tiết sản phẩm nằm trong \tbl{products}.
\end{itemize}
\end{block}
\column{0.48\textwidth}
\begin{block}{JOIN ghép lại thông tin}
\begin{itemize}
  \item Đơn hàng + tên khách hàng.
  \item Chi tiết đơn hàng + tên sản phẩm.
  \item Khách hàng + nhân viên phụ trách.
\end{itemize}
\end{block}
\end{columns}
\vspace{0.5em}
\centering
\texttt{orders.customerNumber = customers.customerNumber}
\end{frame}

% 11
\begin{frame}[fragile]{Cú pháp tổng quát của \texttt{INNER JOIN}}
\begin{lstlisting}
SELECT table1.column1, table2.column2
FROM table1
INNER JOIN table2
    ON table1.key_column = table2.key_column;
\end{lstlisting}
\begin{itemize}
  \item \kw{FROM table1}: bảng bắt đầu truy vấn.
  \item \kw{INNER JOIN table2}: bảng cần nối thêm.
  \item \kw{ON}: điều kiện nối hai bảng.
  \item Chỉ giữ các dòng có dữ liệu khớp ở cả hai bảng.
\end{itemize}
\end{frame}

% 12
\begin{frame}[fragile]{Ví dụ: \texttt{orders} JOIN \texttt{customers}}
\begin{lstlisting}
SELECT
    orders.orderNumber,
    orders.orderDate,
    orders.status,
    customers.customerName
FROM orders
INNER JOIN customers
    ON orders.customerNumber = customers.customerNumber;
\end{lstlisting}
\begin{block}{Ý nghĩa}
Mỗi đơn hàng được bổ sung tên khách hàng tương ứng thông qua \col{customerNumber}.
\end{block}
\end{frame}

% 13
\begin{frame}[fragile]{Đặt bí danh cho bảng}
\begin{lstlisting}
SELECT
    o.orderNumber,
    o.orderDate,
    o.status,
    c.customerName
FROM orders AS o
INNER JOIN customers AS c
    ON o.customerNumber = c.customerNumber;
\end{lstlisting}
\begin{itemize}
  \item Bí danh làm câu lệnh ngắn và dễ đọc.
  \item Có thể bỏ \kw{AS}: \texttt{FROM orders o JOIN customers c ...}
  \item Khi mới học, nên dùng \kw{AS} để rõ nghĩa.
\end{itemize}
\end{frame}

% 14
\begin{frame}[fragile]{JOIN \texttt{products} và \texttt{productlines}}
\begin{lstlisting}
SELECT
    p.productCode,
    p.productName,
    p.productLine,
    pl.textDescription
FROM products AS p
INNER JOIN productlines AS pl
    ON p.productLine = pl.productLine;
\end{lstlisting}
\begin{block}{Cặp khóa}
\centering \texttt{products.productLine = productlines.productLine}
\end{block}
\end{frame}

% 15
\begin{frame}[fragile]{JOIN + \texttt{WHERE} + \texttt{ORDER BY} + \texttt{LIMIT}}
\begin{lstlisting}
SELECT
    p.productCode,
    p.productName,
    p.productLine,
    p.MSRP
FROM products AS p
INNER JOIN productlines AS pl
    ON p.productLine = pl.productLine
WHERE p.productLine = 'Classic Cars'
ORDER BY p.MSRP DESC
LIMIT 10;
\end{lstlisting}
\end{frame}

% 16
\begin{frame}[fragile]{JOIN \texttt{employees} và \texttt{offices}}
\begin{lstlisting}
SELECT
    e.employeeNumber,
    e.lastName,
    e.firstName,
    e.jobTitle,
    o.city,
    o.country
FROM employees AS e
INNER JOIN offices AS o
    ON e.officeCode = o.officeCode;
\end{lstlisting}
\begin{block}{Cặp khóa}
\centering \texttt{employees.officeCode = offices.officeCode}
\end{block}
\end{frame}

% 17
\begin{frame}[fragile]{JOIN \texttt{customers} và \texttt{employees}}
\begin{lstlisting}
SELECT
    c.customerNumber,
    c.customerName,
    c.country,
    e.employeeNumber,
    e.lastName,
    e.firstName
FROM customers AS c
INNER JOIN employees AS e
    ON c.salesRepEmployeeNumber = e.employeeNumber;
\end{lstlisting}
\hint{\texttt{INNER JOIN} chỉ trả về khách hàng đã có nhân viên bán hàng phụ trách.}
\end{frame}

% 18
\begin{frame}[fragile]{JOIN \texttt{orderdetails} và \texttt{products}}
\begin{lstlisting}
SELECT
    od.orderNumber,
    od.productCode,
    p.productName,
    od.quantityOrdered,
    od.priceEach
FROM orderdetails AS od
INNER JOIN products AS p
    ON od.productCode = p.productCode;
\end{lstlisting}
\begin{block}{Cặp khóa}
\centering \texttt{orderdetails.productCode = products.productCode}
\end{block}
\end{frame}

% 19
\begin{frame}[fragile]{Tạo cột tính toán khi JOIN}
\begin{lstlisting}
SELECT
    od.orderNumber,
    p.productName,
    od.quantityOrdered,
    od.priceEach,
    od.quantityOrdered * od.priceEach AS lineTotal
FROM orderdetails AS od
INNER JOIN products AS p
    ON od.productCode = p.productCode;
\end{lstlisting}
\hint{Có thể dùng biểu thức trong \kw{SELECT} để tạo cột mới trong kết quả.}
\end{frame}

% 20
\begin{frame}[fragile]{Nối ba bảng}
\begin{lstlisting}
SELECT
    o.orderNumber,
    o.orderDate,
    o.status,
    c.customerName,
    od.productCode,
    od.quantityOrdered,
    od.priceEach
FROM orders AS o
INNER JOIN customers AS c
    ON o.customerNumber = c.customerNumber
INNER JOIN orderdetails AS od
    ON o.orderNumber = od.orderNumber;
\end{lstlisting}
\end{frame}

% 21
\begin{frame}{Ý nghĩa của truy vấn nối ba bảng}
\begin{enumerate}
  \item Nối \tbl{orders} với \tbl{customers} để biết đơn hàng thuộc khách hàng nào.
  \item Nối tiếp với \tbl{orderdetails} để biết đơn hàng có những dòng sản phẩm nào.
  \item Mỗi lần thêm bảng cần thêm một điều kiện \kw{ON} rõ ràng.
\end{enumerate}
\vspace{0.5em}
\centering
\texttt{orders -> customers}\quad và\quad \texttt{orders -> orderdetails}
\end{frame}

% 22
\begin{frame}[fragile]{Nối bốn bảng: thêm \texttt{products}}
\begin{lstlisting}
SELECT
    o.orderNumber,
    o.orderDate,
    c.customerName,
    p.productName,
    od.quantityOrdered,
    od.priceEach
FROM orders AS o
INNER JOIN customers AS c
    ON o.customerNumber = c.customerNumber
INNER JOIN orderdetails AS od
    ON o.orderNumber = od.orderNumber
INNER JOIN products AS p
    ON od.productCode = p.productCode;
\end{lstlisting}
\end{frame}

% 23
\begin{frame}[fragile]{Nối năm bảng: thêm \texttt{productlines}}
\begin{lstlisting}
SELECT
    o.orderNumber,
    o.orderDate,
    c.customerName,
    p.productName,
    p.productLine,
    pl.textDescription,
    od.quantityOrdered,
    od.priceEach
FROM orders AS o
INNER JOIN customers AS c
    ON o.customerNumber = c.customerNumber
INNER JOIN orderdetails AS od
    ON o.orderNumber = od.orderNumber
INNER JOIN products AS p
    ON od.productCode = p.productCode
INNER JOIN productlines AS pl
    ON p.productLine = pl.productLine;
\end{lstlisting}
\end{frame}

% 24
\begin{frame}{Chuỗi nối nhiều bảng}
\begin{block}{Chuỗi dữ liệu thường dùng}
\centering
\texttt{customers -> orders -> orderdetails -> products -> productlines}
\end{block}
\begin{itemize}
  \item Khi thêm bảng, phải biết bảng mới nối với bảng nào đã xuất hiện trong truy vấn.
  \item Không nhất thiết bảng mới phải nối trực tiếp với bảng đầu tiên trong \kw{FROM}.
  \item Mỗi bảng nên có bí danh ngắn: \texttt{o}, \texttt{c}, \texttt{od}, \texttt{p}, \texttt{pl}.
\end{itemize}
\end{frame}

% 25
\begin{frame}[fragile]{Cú pháp \texttt{LEFT JOIN}}
\begin{lstlisting}
SELECT table1.column1, table2.column2
FROM table1
LEFT JOIN table2
    ON table1.key_column = table2.key_column;
\end{lstlisting}
\begin{itemize}
  \item Giữ lại toàn bộ dòng của bảng bên trái.
  \item Nếu không có dòng khớp ở bảng bên phải, các cột bên phải là \kw{NULL}.
  \item Dùng khi muốn tìm dữ liệu thiếu hoặc giữ bảng chính đầy đủ.
\end{itemize}
\end{frame}

% 26
\begin{frame}[fragile]{Ví dụ \texttt{LEFT JOIN}: khách hàng và nhân viên}
\begin{lstlisting}
SELECT
    c.customerNumber,
    c.customerName,
    c.salesRepEmployeeNumber,
    e.employeeNumber,
    e.lastName,
    e.firstName
FROM customers AS c
LEFT JOIN employees AS e
    ON c.salesRepEmployeeNumber = e.employeeNumber;
\end{lstlisting}
\hint{Khách hàng chưa có nhân viên phụ trách vẫn xuất hiện, nhưng phần nhân viên là \kw{NULL}.}
\end{frame}

% 27
\begin{frame}[fragile]{Tìm dòng không có dữ liệu liên quan}
\begin{lstlisting}
SELECT
    c.customerNumber,
    c.customerName,
    c.salesRepEmployeeNumber
FROM customers AS c
LEFT JOIN employees AS e
    ON c.salesRepEmployeeNumber = e.employeeNumber
WHERE e.employeeNumber IS NULL;
\end{lstlisting}
\begin{itemize}
  \item \kw{LEFT JOIN} giữ toàn bộ khách hàng.
  \item \kw{WHERE e.employeeNumber IS NULL} giữ các dòng không có nhân viên khớp.
\end{itemize}
\end{frame}

% 28
\begin{frame}{So sánh \texttt{INNER JOIN} và \texttt{LEFT JOIN}}
\scriptsize
\begin{tabularx}{\textwidth}{l X X}
\toprule
Tiêu chí & \texttt{INNER JOIN} & \texttt{LEFT JOIN} \\
\midrule
Kết quả trả về & Chỉ dòng khớp ở cả hai bảng & Tất cả dòng bảng trái, kèm dòng khớp nếu có \\
Dòng không có liên hệ & Bị loại khỏi kết quả & Vẫn xuất hiện \\
Cột từ bảng phải nếu không khớp & Không có dòng đó & Có giá trị \texttt{NULL} \\
Dùng khi & Chỉ cần dữ liệu đã có liên hệ & Muốn giữ toàn bộ dữ liệu của bảng chính \\
Ví dụ & Đơn hàng kèm khách hàng & Tất cả khách hàng, kể cả chưa có đơn hàng \\
\bottomrule
\end{tabularx}
\end{frame}

% 29
\begin{frame}[fragile]{\texttt{SELF JOIN}: nối bảng với chính nó}
\begin{lstlisting}
SELECT
    e.employeeNumber AS employeeId,
    e.lastName AS employeeLastName,
    e.firstName AS employeeFirstName,
    m.employeeNumber AS managerId,
    m.lastName AS managerLastName,
    m.firstName AS managerFirstName
FROM employees AS e
LEFT JOIN employees AS m
    ON e.reportsTo = m.employeeNumber;
\end{lstlisting}
\begin{itemize}
  \item \texttt{e}: nhân viên.
  \item \texttt{m}: quản lý trực tiếp.
\end{itemize}
\end{frame}

% 30
\begin{frame}{Khi nào dùng \texttt{SELF JOIN}?}
\begin{block}{Điều kiện}
Một bảng có cột tham chiếu đến khóa chính của chính bảng đó.
\end{block}
\begin{itemize}
  \item Trong \texttt{classicmodels}: \col{employees.reportsTo} tham chiếu đến \col{employees.employeeNumber}.
  \item Dùng hai bí danh khác nhau cho cùng một bảng.
  \item \kw{LEFT JOIN} giúp giữ cả nhân viên không có quản lý trực tiếp.
\end{itemize}
\end{frame}

% 31
\begin{frame}[fragile]{Kết hợp \texttt{JOIN} với các mệnh đề khác}
\begin{lstlisting}
SELECT ...
FROM table1 AS a
JOIN table2 AS b
    ON a.key = b.key
WHERE ...
ORDER BY ...
LIMIT ...;
\end{lstlisting}
\begin{block}{Thứ tự thường viết}
\centering
\kw{SELECT} $\rightarrow$ \kw{FROM/JOIN/ON} $\rightarrow$ \kw{WHERE} $\rightarrow$ \kw{ORDER BY} $\rightarrow$ \kw{LIMIT}
\end{block}
\end{frame}

% 32
\begin{frame}[fragile]{Ví dụ tổng hợp: top 10 dòng đơn hàng}
\begin{lstlisting}
SELECT
    o.orderNumber,
    c.customerName,
    p.productName,
    od.quantityOrdered,
    od.priceEach,
    od.quantityOrdered * od.priceEach AS lineTotal
FROM orders AS o
INNER JOIN customers AS c
    ON o.customerNumber = c.customerNumber
INNER JOIN orderdetails AS od
    ON o.orderNumber = od.orderNumber
INNER JOIN products AS p
    ON od.productCode = p.productCode
ORDER BY lineTotal DESC
LIMIT 10;
\end{lstlisting}
\end{frame}

% 33
\begin{frame}[fragile]{Lỗi 1: Quên điều kiện \texttt{ON}}
\begin{columns}[T,onlytextwidth]
\column{0.48\textwidth}
\begin{block}{Sai}
\begin{lstlisting}
SELECT o.orderNumber, c.customerName
FROM orders AS o
INNER JOIN customers AS c;
\end{lstlisting}
\end{block}
\column{0.48\textwidth}
\begin{block}{Đúng}
\begin{lstlisting}
SELECT o.orderNumber, c.customerName
FROM orders AS o
INNER JOIN customers AS c
    ON o.customerNumber = c.customerNumber;
\end{lstlisting}
\end{block}
\end{columns}
\end{frame}

% 34
\begin{frame}[fragile]{Lỗi 2: Nối sai cột}
\begin{columns}[T,onlytextwidth]
\column{0.48\textwidth}
\begin{block}{Sai}
\begin{lstlisting}
SELECT o.orderNumber, c.customerName
FROM orders AS o
INNER JOIN customers AS c
    ON o.orderNumber = c.customerNumber;
\end{lstlisting}
\end{block}
\column{0.48\textwidth}
\begin{block}{Đúng}
\begin{lstlisting}
SELECT o.orderNumber, c.customerName
FROM orders AS o
INNER JOIN customers AS c
    ON o.customerNumber = c.customerNumber;
\end{lstlisting}
\end{block}
\end{columns}
\end{frame}

% 35
\begin{frame}[fragile]{Lỗi 3: Tên cột bị mơ hồ}
\begin{columns}[T,onlytextwidth]
\column{0.48\textwidth}
\begin{block}{Dễ lỗi / khó đọc}
\begin{lstlisting}
SELECT customerNumber, orderNumber
FROM customers AS c
INNER JOIN orders AS o
    ON c.customerNumber = o.customerNumber;
\end{lstlisting}
\end{block}
\column{0.48\textwidth}
\begin{block}{Rõ ràng}
\begin{lstlisting}
SELECT c.customerNumber, o.orderNumber
FROM customers AS c
INNER JOIN orders AS o
    ON c.customerNumber = o.customerNumber;
\end{lstlisting}
\end{block}
\end{columns}
\end{frame}

% 36
\begin{frame}[fragile]{Lỗi 4: \texttt{WHERE} làm mất ý nghĩa \texttt{LEFT JOIN}}
\begin{lstlisting}
-- May remove customers without orders
SELECT c.customerNumber, c.customerName, o.orderNumber
FROM customers AS c
LEFT JOIN orders AS o
    ON c.customerNumber = o.customerNumber
WHERE o.status = 'Shipped';
\end{lstlisting}
\begin{lstlisting}
-- Keep all customers, filter orders in ON
SELECT c.customerNumber, c.customerName, o.orderNumber, o.status
FROM customers AS c
LEFT JOIN orders AS o
    ON c.customerNumber = o.customerNumber
   AND o.status = 'Shipped';
\end{lstlisting}
\end{frame}

% 37
\begin{frame}{Bài tập thực hành theo nhóm}
\scriptsize
\begin{tabularx}{\textwidth}{l X}
\toprule
Phần & Nội dung luyện tập \\
\midrule
2--8 & CSDL mẫu, môi trường, quan hệ bảng, SELECT/WHERE/ORDER/LIMIT \\
9--10 & Cú pháp INNER JOIN và bí danh bảng \\
11--15 & JOIN hai bảng phổ biến \\
16--18 & Nối 3, 4, 5 bảng \\
19--21 & LEFT JOIN, tìm dòng không liên hệ, so sánh với INNER JOIN \\
22 & SELF JOIN \\
23--25 & Kết hợp mệnh đề, sửa lỗi JOIN, bài tổng hợp \\
\bottomrule
\end{tabularx}
\end{frame}

% 38
\begin{frame}{Chiến lược giải bài JOIN}
\begin{enumerate}
  \item Viết truy vấn đơn giản trước: \texttt{SELECT ... FROM main\_table}.
  \item Thêm từng \kw{JOIN} một, kiểm tra cặp khóa trong \kw{ON}.
  \item Chỉ sau khi nối đúng mới thêm \kw{WHERE}.
  \item Với \kw{LEFT JOIN}, cẩn thận khi lọc cột của bảng bên phải trong \kw{WHERE}.
  \item Dùng bí danh bảng để tránh tên cột mơ hồ.
\end{enumerate}
\end{frame}

% 39
\begin{frame}{Tóm tắt}
\begin{block}{Ý chính}
\begin{itemize}
  \item \kw{INNER JOIN}: chỉ lấy dòng khớp ở cả hai bảng.
  \item \kw{LEFT JOIN}: giữ toàn bộ dòng của bảng trái.
  \item \kw{SELF JOIN}: nối một bảng với chính nó bằng hai bí danh.
  \item \kw{ON}: nơi khai báo quan hệ nối bảng.
  \item \kw{WHERE}: nơi lọc dòng sau khi đã nối.
\end{itemize}
\end{block}
\end{frame}

% 40
\begin{frame}{Câu hỏi kiểm tra nhanh}
\begin{enumerate}
  \item Khi nào nên dùng \kw{LEFT JOIN} thay vì \kw{INNER JOIN}?
  \item Vì sao cần đặt bí danh bảng khi truy vấn có nhiều bảng?
  \item Muốn lấy tên sản phẩm trong chi tiết đơn hàng, cần nối bảng nào?
  \item Vì sao điều kiện lọc ở \kw{WHERE} có thể làm mất ý nghĩa của \kw{LEFT JOIN}?
  \item \kw{SELF JOIN} khác gì so với \kw{INNER JOIN} thông thường?
\end{enumerate}
\end{frame}

\end{document}
